\section*{Chapitre \ref{ch:mx:dark-pools-and-blocks} --- Dark pools et transactions de bloc}

\begin{solution}{exo:mx:dark-pools-and-blocks:1}
Le mid vaut 20.01~: l'acheteur économise $20.02-20.01$ et le vendeur $20.01-20.00$, soit un cent par action chacun, ce qui fait 100~USD sur 10\,000 actions.
\end{solution}

\begin{solution}{exo:mx:dark-pools-and-blocks:2}
Rien avec la vente de 300 actions (sous le minimum)~; 1\,500 actions avec la seconde. 3\,500 actions restent en attente.
\end{solution}

\begin{solution}{exo:mx:dark-pools-and-blocks:3}
$28/40=70~\%$. Chaque invitation arrivée à échéance signalait au participant qu'une grande contrepartie existait sans qu'elle soit exécutée~: un taux faible trahit un participant qui utilise
des \omterm{def:mx:dark-pools-and-blocks:conditional}{ordres conditionnels} pour sonder l'information, et la plateforme le restreint pour protéger les autres.
\end{solution}

\begin{solution}{exo:mx:dark-pools-and-blocks:4}
La transaction fait monter $P^\ast$ de $0.0005\times500=0.25$ tick, soit 0.25 point de base à 100.00, sur les 15\,000 actions restant à acheter~: $0.25\times15\,000/20\,000=0.1875$
points de base de l'ordre.
\end{solution}

\begin{solution}{exo:mx:dark-pools-and-blocks:5}
Les \omterm{def:mx:why-there-is-a-spread:informed}{traders informés} négocient lorsque $P^\ast$ s'écarte du mid et tous du même côté~; la liquidité dark dont ils ont besoin, du côté opposé, a été postée par
des traders non informés patients, placés au hasard, et les informés l'épuisent. Les preneurs de liquidité non informés arrivent des deux côtés et la trouvent plus souvent.
\end{solution}

\begin{solution}{exo:mx:dark-pools-and-blocks:6}
L'ordre enfant $i$ de $q=Q/n$ est suivi de $Q-iq$ actions qui paient son impact $\lambda q$~: $\sum_{i=1}^{n}\lambda q(Q-iq)/Q=\lambda q\,(n-(n+1)/2)=\lambda Q/2\,(1-1/n)$.
Avec $\lambda=0.0005$, $Q=20\,000$, $n=100$~: 4.95 ticks, soit 4.95 points de base à 100.00.
\end{solution}

\begin{solution}{exo:mx:dark-pools-and-blocks:7}
Un minimum de 200 bloque toujours les sondes de 100 actions de l'\omterm{def:mx:dark-pools-and-blocks:ping}{ordre sonde}~: aucune fuite, l'ordre entier s'exécute dans le dark pool, en 1\,020 secondes environ en moyenne, pour un
shortfall de 7.0 points de base (erreur type 0.9) contre 12.5 pour l'\omterm{def:mx:order-types-and-their-uses:peg}{ordre indexé au mid} non protégé et 9.5 pour un minimum de 1\,000~; la dérive est de 13.6. Un minimum juste
au-dessus de la taille des sondes conserve la protection et perd bien moins de flux naturel.
\end{solution}

\begin{solution}{exo:mx:dark-pools-and-blocks:8}
Les exécutions au mid font économiser la fourchette, mais pas le reste~: le cours a bougé pendant le travail de l'ordre, parce que ses exécutions ont prélevé du flux sur le marché lit et
parce que d'autres en ont eu connaissance. Mesurer le shortfall par rapport au prix d'arrivée, la dérive avant la fin de l'exécution et les mark-out des contreparties~; dans la
simulation, l'ordre dark non protégé a coûté 12.5 points de base, tout au mid.
\end{solution}

\begin{solution}{pb:mx:dark-pools-and-blocks:1}
\textbf{1.} La moitié de la fourchette pour chaque côté~; le meilleur cours acheteur et le meilleur cours vendeur du marché lit.
\textbf{2.} Une copie à jour des cotations du marché lit~: contrôle N, envoyé après chaque changement du sommet du carnet du marché lit.
\textbf{3.} La \omterm{def:mx:dark-pools-and-blocks:minqty}{quantité minimale d'exécution} qu'accepte un ordre dark~; elle tient à l'écart les petits \omterm{def:mx:dark-pools-and-blocks:ping}{ordres sondes} et réduit le nombre d'exécutions.
\textbf{4.} Une \omterm{def:mx:dark-pools-and-blocks:ioi}{indication d'intérêt} (IOI) non contraignante et au repos, qui doit être confirmée lorsqu'elle est invitée~; le \omterm{def:mx:dark-pools-and-blocks:firmup}{taux de confirmation} mesure la part des invitations qu'un participant confirme.
\textbf{5.} Les deux côtés apprennent que la contrepartie existe, et l'un d'eux n'a rien payé pour l'apprendre.
\textbf{6.} Les \omterm{def:mx:why-there-is-a-spread:informed}{traders informés} se regroupent d'un côté et font face à un risque d'exécution plus élevé dans le dark pool~; ils préfèrent donc la bourse, tandis que les non-informés préfèrent le dark pool.
\textbf{7.} 6.3~\% et 17.5~\%.
\textbf{8.} Dark~: aucune capture, $+0.11$ cent de variation~; fournisseurs de liquidité sur le marché lit~: 0.51 cent de capture, $-0.13$ de variation.
\textbf{9.} Rien ne montre que les transactions de bloc dans le dark pool entravent la \omterm{def:mx:fragmentation-and-routing:discovery}{formation des prix}.
\textbf{10.} L'\omterm{def:mx:dark-pools-and-blocks:ping}{ordre sonde} émet une sonde de 100 actions toutes les deux secondes et négocie 500 actions par anticipation après deux exécutions consécutives~; 30~\% des vendeurs de bloc ne confirment jamais et achètent
2\,000 actions par anticipation.
\textbf{11.} Voir l'exercice~6~: 4.95 points de base.
\textbf{12.} $\lambda\times$ chaque transaction déclenchée par anticipation $\times$ les actions du fonds restant à exécuter, rapporté à la taille de l'ordre.
\textbf{13.} Marché lit~: 3.8, 7.3, 4.95 d'impact propre~; dark pool~: 12.5, 26.1, 10.9 de fuite~; dark pool avec minimum~: 9.5, 11.9, 1.4 d'impact propre~; bloc~: $-2.9$, $-1.7$, 0.6
fuite (erreurs types de 0.5 à 2.3).
\textbf{14.} L'\omterm{def:mx:dark-pools-and-blocks:ping}{ordre sonde} l'a détecté 88 fois et a acheté 44\,000 actions par anticipation~: la fuite (10.9 points de base) est supérieure à l'impact propre de l'ordre sur le marché lit (4.95).
\textbf{15.} Il arrête la fuite~; seuls 46~\% s'exécutent dans le dark pool et le reste est acheté à la fin.
\textbf{16.} Leurs vendeurs apportent un volume qui, sinon, n'aurait pas été négocié~; ce n'est plus le cas avec des contreparties qui fuient ou sont informées, ou lorsque les vendeurs du bloc auraient
vendu sur le carnet du marché lit.
\textbf{17.} ITG~: un desk non divulgué a négocié contre les abonnés de son dark pool en utilisant leurs ordres en temps réel (20.3~millions de dollars). Barclays~: a présenté de façon trompeuse la manière dont
elle surveillait son dark pool (70~millions de dollars). Credit Suisse~: a présenté de façon trompeuse son système de notation, a accepté des ordres à moins d'un centime et a signalé aux sociétés de trading haute fréquence les
ordres de ses clients (84.3~millions de dollars).
\textbf{18.} Le double plafond de volume a été remplacé par un plafond de volume unique pour la dérogation au prix de référence, à compter de 18~mois après l'entrée en
vigueur de la révision de MiFIR.
\textbf{19.} \emph{Résultat nommé}~: sur la fenêtre de 30 minutes, le mid a dérivé de 26.1 points de base contre l'ordre dark non protégé et de 7.3 contre
l'ordre sur le marché lit, et la fuite comptabilisée de l'ordre dark a coûté 10.9 points de base~; une quantité minimale supérieure à la taille des sondes l'a supprimée.
\textbf{20.} Cibler d'abord les plateformes de blocs où se trouvent des vendeurs naturels, et ne recourir au dark pool que derrière une quantité minimale supérieure à toute sonde.
\end{solution}

\begin{solution}{iq:mx:dark-pools-and-blocks:1}
Une plateforme qui n'affiche aucun ordre et croise en général au mid du marché lit~: une institution s'en sert pour négocier un volume sans le montrer, en économisant la fourchette et
espérant éviter l'impact.
\iqlookfor{Mid~; absence d'affichage~; volume et fuite.}
\end{solution}

\begin{solution}{iq:mx:dark-pools-and-blocks:2}
Une \omterm{def:mx:dark-pools-and-blocks:minqty}{quantité minimale d'exécution} supérieure aux tailles usuelles des sondes, une répartition aléatoire et partielle entre les plateformes, en évitant celles aux contrôles des participants insuffisants,
les \omterm{def:mx:dark-pools-and-blocks:conditional}{ordres conditionnels} uniquement sur les plateformes qui surveillent les \omterm{def:mx:dark-pools-and-blocks:firmup}{taux de confirmation}, et la surveillance des exécutions pour repérer les signes de sondage.
\iqlookfor{Quantité minimale~; choix des plateformes~; surveillance.}
\end{solution}

\begin{solution}{iq:mx:dark-pools-and-blocks:3}
Théorie (Zhu)~: les dark pools attirent le flux non informé et peuvent concentrer l'information sur la bourse. Résultats empiriques (Comerton-Forde et Putnins)~: les faibles niveaux de
transactions hors bloc en dark pool sont bénignes, les niveaux élevés nuisibles, et les blocs ne le sont pas.
\iqlookfor{Argument de sélection~; non-linéarité des résultats empiriques.}
\end{solution}

\begin{solution}{iq:mx:dark-pools-and-blocks:4}
Comparer la dérive du cours pendant le travail des ordres à un indice de référence d'ordres similaires, les mark-out des contreparties par plateforme et les mouvements de cours après
les IOI et les invitations~; randomiser le routage lorsque c'est possible pour séparer la fuite de la sélection.
\iqlookfor{Dérive~; mark-out des contreparties~; groupe de comparaison.}
\end{solution}

\begin{solution}{iq:mx:dark-pools-and-blocks:5}
Par ordre~: propriétaire, sens, taille, minimum, heure, et s'il figure dans une invitation ouverte. Par invitation~: les deux ordres, la taille, l'heure de début, qui a confirmé. Par
participant~: invitations et confirmations (le taux), et toute restriction. Appariement par priorité temporelle entre les ordres libres dont les tailles atteignent les deux minimums.
\iqlookfor{État de l'invitation~; minuterie~; notation des participants.}
\end{solution}

\begin{solution}{iq:mx:dark-pools-and-blocks:6}
Le mid économise la fourchette, mais l'ordre a fui ou a prélevé du flux sur le marché lit, et le cours a dérivé pendant son travail~; mesurer la dérive et
les mark-out des contreparties, ajouter une quantité minimale et comparer avec les exécutions en marché lit et en bloc d'ordres similaires.
\iqlookfor{Fuite et impact au-delà de la fourchette~; comment les mesurer.}
\end{solution}
